%  LaTeX support: latex@mdpi.com 
%  For support, please attach all files needed for compiling as well as the log file, and specify your operating system, LaTeX version, and LaTeX editor.

%=================================================================
\documentclass[coasts,article,submit,pdftex,moreauthors]{Definitions/mdpi} 

%=================================================================

% MDPI internal commands
\firstpage{1} 
\makeatletter 
\setcounter{page}{\@firstpage} 
\makeatother
\pubvolume{1}
\issuenum{1}
\articlenumber{0}
\pubyear{2024}
\copyrightyear{2024}
%\externaleditor{Academic Editor: Firstname Lastname}
\datereceived{25 December 2023} 
%\daterevised{} % Only for the journal Acoustics
\dateaccepted{} 
\datepublished{} 
\hreflink{https://doi.org/} % If needed use \linebreak
%\doinum{}

%=================================================================
% Add packages and commands here. The following packages are loaded in our class file: fontenc, inputenc, calc, indentfirst, fancyhdr, graphicx, epstopdf, lastpage, ifthen, lineno, float, amsmath, setspace, enumitem, mathpazo, booktabs, titlesec, etoolbox, tabto, xcolor, soul, multirow, microtype, tikz, totcount, changepage, attrib, upgreek, cleveref, amsthm, hyphenat, natbib, hyperref, footmisc, url, geometry, newfloat, caption

\graphicspath{{figures/}} % path to folder with figures
\usepackage{subcaption}
\usepackage{listings}
\usepackage{xcolor}
\usepackage{gensymb} % for \textdegree
\usepackage{tabularx}

%=================================================================
% Full title of the paper (Capitalized)
\Title{Random Forest Classifier Algorithm of GRASS GIS for Satellite Image Processing in the Bight of Sofala, Eastern Mozambique}

% MDPI internal command: Title for citation in the left column
\TitleCitation{Random Forest Classifier Algorithm of GRASS GIS for Satellite Image Processing in the Bight of Sofala, Eastern Mozambique}

% Author Orchid ID: enter ID or remove command
\newcommand{\orcidauthorA}{0000-0002-5759-1089} % Polina Add \orcidA{} behind the author's name

% Authors, for the paper (add full first names)
\Author{Polina Lemenkova $^{1}$*\orcidA{}}

%\longauthorlist{yes}

% MDPI internal command: Authors, for metadata in PDF
\AuthorNames{Polina Lemenkova}

% MDPI internal command: Authors, for citation in the left column
\AuthorCitation{Lemenkova, P.}
% If this is a Chicago style journal: Lastname, Firstname, Firstname Lastname, and Firstname Lastname.

% Affiliations / Addresses (Add [1] after \address if there is only one affiliation.)
\address{%
$^{1}$ \quad Department of Geoinformatics, Faculty of Digital and Analytical Sciences, Universität Salzburg, Schillerstraße 30, A-5020 Salzburg, Austria.
}

% Contact information of the corresponding author
\corres{Correspondence: polina.lemenkova@plus.ac.at; Tel.: +43-677-6173-2772 (P.L.)}

% Current address and/or shared authorship
%\firstnote{Current address: Affiliation 3.} 

% Abstract (Do not insert blank lines, i.e. \\) 
\abstract{Mapping coastal regions is important for environmental assessment and monitoring spatio-temporal changes. Although the traditional cartographic methods using Geographic Information System (GIS) is applicable for image classification, the Machine Learning (ML) methods present more advantageous solutions of pattern-finding tasks such as automated detection of landscape patches in heterogeneous landscapes. This study intended to discriminate the landscape patterns along the eastern coasts of Mozambique using ML modules of Geographic Resources Analysis Support System (GRASS) GIS. The Random Forest (RF) algorithm of module 'r.learn.train' has been used to map the coastal landscapes of eastern shoreline of the Bight of Sofala, using remote sensing (RS) data at multiple temporal scales. The dataset included a time series Landsat 8-9 OLI/TIRS imagery collected at dry period during 2015, 2018 and 2023. The supervised sub-pixel classification of RS rasters was supported by the Scikit-Learn ML package of Python embedded in GRASS GIS. The Bight of Sofala is characterised by the distributed marine ecosystems with dominated swamp wetlands and mangrove forests located in the mixed saline-fresh waters along the eastern coast of Mozambique in the Beira estuary entering the Mozambique Channel. To evaluate short-term temporal dynamics in this region, advanced methods of cartographic scripting using ML modules of GRASS GIS were applied. The methodology included the training data extraction from the raster images, supervised ML using ensemble classification and regression tree method by RF and cross-validation tasks. Technically, the paper demonstrated the advantages of using the ML algorithm for RS data classification for environmental monitoring of coastal areas. Specifically, the high robustness, versatility and applicability of the RF algorithm is useful for processing satellite images in the shoreline regions with sensitive environmental parameters. Due to sensitivity of RF to noise and outliers in the pixel matrix, its application to image processing is beneficial for reducing overfitting by averaging multiple decision trees. This resulted in reduced data processing time and increased accuracy of Landsat 8-9 OLI/TIRS imagery. The integration of Earth Observation data, processed using multiple decision tree classifier using ML methods and land cover characteristics enabled to detect recent changes in the coastal ecosystem of Mozambique, East Africa.
}

% Keywords
\keyword{Machine learning; ensemble learning; GRASS GIS; Scikit-Learn; Python; random forest; satellite image; image processing; Africa; image classification} 

% The fields PACS, MSC, and JEL may be left empty or commented out if not applicable
%\PACS{J0101}
%\MSC{}
%\JEL{}

\begin{document}

%----------- section ----------->
\section{Introduction}

%----------- section ----------->
\section{Materials and Methods}

The data processing was performed on the MacBook Apple Sonoma (arm64). The Cooperative Computing Tools (CCTools) which enable large scale distributed computations from clusters were updated and installed using sudo as follows: 'sudo port clean cctools' and 'sudo port -v install cctools'. After the CCTools and XCode were installed and activated, necessary updates in the libSystem of MacPorts were performed including the deactivated library libunwind-headers: 'sudo port -f deactivate libunwind-headers'. The installation of the latest version of GRASS GIS 8.3.1 used in this research was performed using 'sudo port install grass'. The compilation of the GRASS GIS 8.3 included the dependancies on the essential Python's libraries and requirements of ML packages. The toolchain GNU Compiler Collection (GCC) 13 was installed accordingly for code compilation, supporting library dependencies through linking and conversion of GRASS GIS scripts into the binary executable format to assembly for executable files. The wxPython was installed additionally using the 'sudo port install py311-wxpython-4.0' call from the MacPorts.

%----------- subsection ----------->
\subsection{Subsection}

%----------- subsection ----------->
\subsection{Subsection}

%----------- subsection ----------->
\subsection{Subsection}

%----------- section ----------->
\section{Results}

%----------- section ----------->
\section{Discussion}

%----------- section ----------->
\section{Conclusions}

\funding{The publication was funded by the Editorial Office of Coasts, Multidisciplinary Digital Publishing Institute (MDPI), by providing 100\% discount for the APC of this manuscript.}

\institutionalreview{Not applicable.}

\informedconsent{Not applicable.}

\dataavailability{Not applicable} 

\acknowledgments{The authors thank the reviewers for reading and review of this manuscript.}

\conflictsofinterest{The authors declare no conflict of interest.} 

\abbreviations{Abbreviations}{
The following abbreviations are used in this manuscript:\\

\noindent 
\begin{tabular}{@{}ll}
CCTools & Cooperative Computing Tools \\
GRASS & Geographic Resources Analysis Support System \\
GIS & Geographic Information System \\
GCC & GNU Compiler Collection \\
Landsat OLI/TIRS & Landsat Operational Land Imager and Thermal Infrared Sensor \\
ML & Machine. Learning \\
RF & Random Forest \\
\end{tabular}
}

%%%%%%%%%%%%%%%%%%%%%%%%%%%%%%%%%%%%%%%%%%
%% Optional
\appendixtitles{no} % Leave argument "no" if all appendix headings stay EMPTY (then no dot is printed after "Appendix A"). If the appendix sections contain a heading then change the argument to "yes".
\appendixstart
\appendix
\section[\appendixname~\thesection]{}
\subsection[\appendixname~\thesubsection]{}
The appendix 

\section[\appendixname~\thesection]{}
All appendix sections must be cited in the main text. In the appendices, Figures, Tables, etc. should be labeled, starting with ``A''---e.g., Figure A1, Figure A2, etc.

%%%%%%%%%%%%%%%%%%%%%%%%%%%%%%%%%%%%%%%%%%
\begin{adjustwidth}{-\extralength}{0cm}
%\printendnotes[custom] % Un-comment to print a list of endnotes

\reftitle{References}
\nocite{*}

%=====================================
% References, variant A: external bibliography
%=====================================
\bibliography{Bibliography.bib}

%%%%%%%%%%%%%%%%%%%%%%%%%%%%%%%%%%%%%%%%%%
\end{adjustwidth}
\end{document}
